\section{Discussion}\label{sec:discussion}

\paragraph{The linear limit is not the attainable depth}
With a clipped actuator the linear stability limit $\alim(g)$ is an upper bound that a cut reaches only if it enters gently and is not disturbed. Two consequences follow for practice. First, the way the tool enters the cut is part of the design: a step entry from rest, an idealised full engagement at $t=0$, saturates the actuator for a few tooth periods and, at 157 of the 251 speeds of the case study, ends in chatter below $0.97\,\alim$, whereas a face entry or a feed ramp reaches the limit (Section~\ref{sec:bistability}). Second, the authority that the feedback spends on the forced vibration must be budgeted explicitly, and with a margin. At \SI{2500}{\rpm} and $g=\SI{1500}{\volt\second\per\metre}$ the forced command at the limit leaves only \SI{0.87}{\volt} free, and \SI{0.73}{\micro\metre} of eccentric runout brings it to the clip at some depth up to $0.9\,\alim$; for DVF designs that keep \SI{10}{\percent} of the voltage free, 1.1--\SI{4.5}{\micro\metre} of runout does the same (Section~\ref{sec:delayed-robust}). A voltage guard for runout and noise therefore belongs in the authority condition of Section~\ref{sec:feasibility}. The certificates of Section~\ref{sec:certificates} bound from below, on collocation models, the instantaneous perturbation that the remaining authority can absorb; persistent disturbances, which would call for invariant-set or $\mathcal L_2$ formulations \citep{hu2002disturbance,fang2004l2}, are not covered.
% src: results/c1_df.npz (V_max - b_max at a_lim, 2500 rpm: 0.87 V); verification/paper/revision_runs.json (runout_to_clip_g1500: 0.73 um at 2500 rpm); results/c4_summary.json (robust: DVF runout tolerance 1.12-4.50 um with a 10 % guard band)

\paragraph{How conservative the certificates are, and why}
Where the simulated basin is finite, the certified kicks are 1.5 to 28 times smaller than the simulated critical kicks, and about 2.3 times smaller near the limit; where no kick up to \SI{3}{\metre\per\second} caused chatter, they are smaller by a factor of at least 240. The critical kick of the collocation model itself is 0.62--0.99 times the simulated one, and averaging the time-periodic coefficient changes it by only $-3$ to $+5\,\%$ for this helical cutter, so the gap lies in the certificate. Re-solving with a looser validity bound locates part of it: with $\delta=1$ instead of 0.5 the certified kick rises by \SI{58}{\percent} at \SI{1090}{\rpm} and $0.15\,\alim$ and by \SI{15}{\percent} at \SI{2870}{\rpm} and $0.3\,\alim$, but by only \SI{0.4}{\percent} at \SI{2500}{\rpm} and $0.9\,\alim$, so the plateau at shallow depths is set largely by the validity bound, and near the limit by the ellipsoidal shape and the sector bound. Near the limit the describing-function estimate of Section~\ref{sec:bistab-df} is within 0.88--1.13 of the simulated basin at almost no cost: the certificate gives a guarantee for the model, the describing function a design estimate. Non-quadratic, piecewise-quadratic or Lur'e-type Lyapunov functions \citep{hu2004lure,hu2006nonquadratic,dai2009piecewise}, and Lyapunov--Krasovskii functionals constructed for the delay equation itself \citep{papachristodoulou2009sos}, are the natural next steps; a certified step entry at productive depths would need a phase-dependent entry condition.
% src: results/c3_summary.json (errata.c2_validate 1.48-27.6; basin model_over_td 0.62-0.99, averaging 0.973-1.052); results/c2_validate.npz (19 points without chatter: certificates 1.67-12.51 mm/s, >= 240 times below 3 m/s); verification/paper/revision_runs.json (beta_vs_delta); results/c1_summary.json (td_over_df at 0.95-0.98 a_lim 0.88-1.13)

\paragraph{What the delayed term buys and what it costs}
Because $\dot x(t)-\dot x(t-T)$ vanishes on the forced orbit, the delayed term adds damping of the chatter deviation without adding to the forced command. Under the same authority rule as DVF it raises the linear limit by a median factor of 1.78 for $g_2=\SI{2000}{\volt\second\per\metre}$, but the new limit exceeds twice the cutter diameter at 222 of 251 speeds (71 for DVF), and with the usable depth capped at \SI{32}{\milli\metre} the median gain is 1.31. At the same forced command the simulated basin widens by a median factor of 1.14 over 22 depth--speed pairs, and up to 2.6 times where the DVF certificate almost vanishes; the step-entry limit rises by a median factor of only 1.03, so the bistable band widens. The price is robustness: the phase of $1-e^{-sT}$ alternates between lead and lag between tooth-passing harmonics, and where it lags it adds to the loop delay. With the assumed second modes, $g_2=\SI{1000}{\volt\second\per\metre}$ tolerates about \SI{200}{\micro\second} of loop delay and $g_2\ge\SI{2000}{\volt\second\per\metre}$ about \SI{100}{\micro\second}, against \SI{300}{\micro\second} for DVF; the delay must follow the spindle speed within about \SI{1}{\percent}, and the term makes the forced command more sensitive to runout. A gain $g_2\approx\SI{1000}{\volt\second\per\metre}$ is a reasonable compromise. A low-pass filter in the delay path, as in repetitive control \citep{hara1988repetitive}, or DVF with notch filters at the tooth-passing harmonics \citep{herzog1996unbalance} might keep part of the benefit with less phase lag; neither was evaluated.
% src: results/c4_summary.json (limits, entry median 1.03, validate2 same_g1_over_dvf median 1.14, robust); verification/paper/revision_runs.json (depth_2D: 222 and 71 speeds above 32 mm; capped median 1.31)

\paragraph{Limitations}
The results are predictions of a model that classifies 17 of the 19 published cutting tests of \citet{ozsoy2026robust} as measured; the two exceptions are controlled cuts that chattered at 0.76 and 0.80 of the predicted limit, and at both the clipped model settles after either entry; the bistability studied here cannot explain the cut at \SI{2800}{\rpm}, and explains the one at \SI{1300}{\rpm} only if the test contained a perturbation of about \SI{25}{\milli\metre\per\second} (Section~\ref{sec:model-case}). The model has one identified mode, uniform over the depth, a rigid tool, a cutting force linear in the chip thickness without edge forces or process damping \citep{altintas2012book,budak2010process}, and an actuator whose only nonlinearity is a hard clip on an exact velocity, without stroke or current limits, amplifier dynamics or sensor noise; the second modes of Section~\ref{sec:delayed-robust} are assumed. The attribution of the bistability to the clip holds for this linear force law, since a nonlinear force law can itself make the cut subcritical \citep{kalmarnagy2001subcritical,dombovari2008estimates}. Outside the basin the model predicts clipped chatter of a few tenths of a millimetre, or divergence, beyond its range of validity. Edge forces would add to the static force and hence to the forced command, so the free voltage is likely overestimated. The certificates hold for collocation models, whose certified size changes by up to about \SI{2}{\percent} when the number of nodes is doubled (\ref{app:lmi}), and for instantaneous perturbations only.
% src: results/cuts_main.json, verification/targets/experimental_points.csv; verification/paper/revision_runs.json (published_cuts, beta_vs_N); README 7, 11.2, 11.7

\paragraph{Towards experimental validation}
The predictions could be tested in two stages. On a hardware-in-the-loop rig, where the real actuator, amplifier and controller act on the identified structure while the regenerative force is emulated \citep{sahu2021bistability,sahu2023adaptive,mancisidor2019mechatronics}, command pulses applied through the actuator at controlled phases, with and without the clip, would measure the critical pulse against the depth directly and repeatably. Cutting tests at two or three speeds where the predicted band is wide would then compare a face entry with an entry into a pre-machined pocket, or a feed hold and restart inside the cut, at depths inside the band, with the runout and the settled command measured and enough repetitions to estimate probabilities, since uncontrolled disturbances make the outcome random.
